\chapter{Vaughan Williams, Ralph} % (fold)
\label{cha:vaughan williams}

\href{https://en.wikipedia.org/wiki/List_of_compositions_by_Ralph_Vaughan_Williams}{List of works} on Wikipedia.


\section{Incidental}

\noindent{The Wasps (1909)} \\
\noindent{The Wasps, Aristophanic Suite (1912)}   \\
\noindent{The Bacchae (1911)}  \\
\noindent{The Death of Tintagiles (1913)} \\
\noindent{The First Nowell, nativity play adapted from medieval pageants by Simona Pakenham; score completed by Roy Douglas (1958)}  \\

\section{Ballet}

\noindent{Old King Cole (1923)} \\
\noindent{On Christmas Night (1926)} \\
\noindent{Job: A Masque for Dancing (1930)} \\
\noindent{The Voice out of the Whirlwind, Motet for mixed choir and organ or orchestra; adapted from "Galliard of the Sons of the Morning" from Job}  \\
\noindent{The Running Set (1933)}  \\
\noindent{The Bridal Day (1938–39)}  \\
\noindent{Epithalamion (1957), Cantata for baritone, chorus and small orchestra}  \\

\section{Orchestral} % (fold)

\noindent{Fantasia for Piano and Orchestra (1896)} \\
\noindent{Serenade in A minor (1898)} \\
\noindent{Heroic Elegy and Triumphal Epilogue (1900)} \\
\noindent{Bucolic Suite (1901)} \\
\noindent{Burley Heath, impression for orchestra (1902–03)} \\
\noindent{The Solent, impression for orchestra (1902–03)} \\
\noindent{Symphony No. 1 'A Sea Symphony' (1903–1909) (with chorus, on texts by Whitman)} \textit{BBC SO, Martyn Brabbinbs}\marginnote{max} \\
\noindent{In the Fen Country, for orchestra (1904)} \\
\noindent{Norfolk Rhapsody No. 1, for orchestra (1906, rev. 1914)} \\
\noindent{Norfolk Rhapsody No. 2, for orchestra (1906, subsequently withdrawn; reconstructed and recorded in 2002)} \\
\noindent{Harnham Down, impression for orchestra (1904–07)} \\
\noindent{Fantasia on a Theme by Thomas Tallis (1910, rev. 1913 and 1919)} \\
\noindent{Symphony No. 2 'A London Symphony' (1911–13; revised 1918, 1920 and 1933)} \textit{BBC SO, Martyn Brabbinbs}\marginnote{high} \\
\noindent{The Lark Ascending for Violin and Orchestra (1914)} \textit{Isabelle van Keulen, NDR Radiophilharmonie}\marginnote{high} \\
\noindent{Symphony No. 3 'Pastoral Symphony' (1921)} \textit{BBC SO, Martyn Brabbinbs}\marginnote{high} \\
\noindent{Concerto Accademico for Violin and String Orchestra (1924–25)} \\
\noindent{Flos Campi for Viola, Wordless Chorus, and Small Orchestra (1925)} \\
\noindent{Piano Concerto in C Major (1926–31)} \\
\noindent{Fantasia on Sussex Folk Tunes for Cello and Orchestra (1929); withdrawn by the composer} \\
\noindent{Symphony No. 4 in F minor (1931–34)} \textit{BBC SO, Martyn Brabbinbs}\marginnote{high} \\
\noindent{Fantasia on "Greensleeves" (1934) (for string orchestra and harp; arranged by Ralph Greaves from Vaughan Williams's treatment of folk tunes in his opera \textit{Sir John in Love})} \\
\noindent{Suite for Viola and Small Orchestra (1934)} \\
\noindent{Two Hymn Tune Preludes (1936) for small orchestra: 1. Eventide; 2. Dominus regit me} \\
\noindent{Symphony No. 5 in D major (1938–43)} \textit{BBC SO, Martyn Brabbinbs}\marginnote{max} \\
\noindent{Five Variants of Dives and Lazarus (1939) for strings and harp} \\
\noindent{Sketches for Cello Concerto (1942–43); incomplete, with 2nd movement completed by David Matthews (2009) as Dark Pastoral} \\
\noindent{Oboe Concerto in A Minor for Oboe and Strings (1944)} \\
\noindent{Symphony No. 6 in E minor (1944–47, rev. 1950)} \textit{BBC SO, Martyn Brabbinbs}\marginnote{02/2022, high} \\
\noindent{Concerto for Two Pianos and Orchestra (1946), arranged by Joseph Cooper in collaboration with the composer} \\
\noindent{Concerto for Two Pianos and Orchestra (1946)} \\
\noindent{Partita for Double String Orchestra (1948), rewritten from Double Trio for string sextet with new finale} \\
\noindent{Fantasia (quasi variazione) on the Old 104th Psalm Tune for Piano, Chorus, and Orchestra (1949)} \\
\noindent{Symphony No. 7 'Sinfonia antartica' (1949–52) (partly based on his music for the film \textit{Scott of the Antarctic})} \textit{Elizabeth Watts, BBC SO, Martyn Brabbinbs}\marginnote{max} \\
\noindent{Concerto Grosso, for three groups of strings, each requiring different levels of technical skill (1950)} \\
\noindent{Romance in D-flat Major for Harmonica and Orchestra (1951), written for Larry Adler} \\
\noindent{Symphony No. 8 in D minor (1953–55)} \textit{BBC SO, Martyn Brabbinbs}\marginnote{02/2022, high} \\
\noindent{Tuba Concerto in F Minor (1954)} \\
\noindent{Symphony No. 9 in E minor (1956–57)} \textit{BBC SO, Martyn Brabbinbs}\marginnote{max} \\
\noindent{Flourish for Glorious John (1957)} \\





